\documentclass[11pt,a4paper,ja=standard]{bxjsarticle}

% ---------- LuaLaTeX用日本語設定 ----------
\usepackage{luatexja}
\usepackage{luatexja-fontspec}
\setmainjfont{Yu Mincho}
\setsansjfont{Yu Gothic}

% ---------- 数学・図式パッケージ ----------
\usepackage{amsmath,amssymb,amsthm}
\usepackage{mathtools}
\usepackage{tikz-cd}
\usepackage{tikz}
\usetikzlibrary{positioning,arrows,calc}

% ---------- レイアウト ----------
\usepackage{geometry}
\geometry{margin=25mm}
\usepackage{hyperref}
\usepackage{enumitem}
\usepackage{xcolor}

% ---------- 定理環境 ----------
\theoremstyle{definition}
\newtheorem{definition}{定義}[section]
\newtheorem{theorem}[definition]{定理}
\newtheorem{lemma}[definition]{補題}
\newtheorem{proposition}[definition]{命題}
\newtheorem{example}[definition]{例}
\newtheorem{remark}[definition]{注意}

% ---------- カスタムマクロ ----------
\newcommand{\layer}{\operatorname{layer}}
\newcommand{\Hom}{\operatorname{Hom}}
\newcommand{\Obj}{\operatorname{Obj}}
\newcommand{\Type}{\mathbf{Type}}
\newcommand{\RankHomLe}{\mathtt{RankHomLe}}
\newcommand{\RankHomEq}{\mathtt{RankHomEq}}
\newcommand{\RankHomD}{\mathtt{RankHomD}}
\newcommand{\RankCatLe}{\mathtt{RankCatLe}}
\newcommand{\Ranked}{\mathtt{Ranked}}

% ---------- タイトル情報 ----------
\title{ランク付き構造の圏論的枠組み：\\
3つのモルフィズム層とそのブリッジ}

\author{%
  TeXファイル生成: Claude Code\\
  Leanコード生成: OpenAI Codex\\
  \medskip\\
  \small 本文書はLean 4による形式化プロジェクトの数学的解説です
}

\date{\today}

% ---------- ドキュメント開始 ----------
\begin{document}

\maketitle

\begin{abstract}
本稿では、ランク関数を持つ構造の圏論的枠組みを構築する。
主な貢献は以下の通りである：
(1) データ層モルフィズム $\RankHomLe$ の定義と圏構造 $\RankCatLe$、
(2) 3つの強度レベル（等式層 $\RankHomEq$、不等式層 $\RankHomLe$、存在層 $\RankHomD$）の定義、
(3) これらの層の間の自然な埋め込み（ブリッジ）の構成。
本研究はLean 4による形式的証明を伴い、圏の公理が構成的に検証されている。
\end{abstract}

\tableofcontents
\clearpage

% ============================================================
\section{序論}
% ============================================================

\subsection{背景と動機}

ランク付き構造は、複雑度理論、計算理論、順序構造の研究など、数学と計算機科学の様々な分野で自然に現れる。
各要素にランク（順序付き集合の値）を割り当てることで、構造の「階層性」や「複雑さ」を捉えることができる。

本研究の目的は、このようなランク付き構造の間の射を圏論的に整理し、
異なる強度の条件（等式、不等式、存在条件）を持つ3つの「層」を定義し、
それらの間の関係を明確化することである。

\subsection{形式化の方針}

本稿で述べる数学的内容は、すべてLean 4による形式的証明を伴っている。
Leanコードの生成にはOpenAI Codexを使用し、TeXファイルの生成にはClaude Codeを使用した。
形式化により、以下の利点が得られる：

\begin{itemize}
\item 圏の公理（結合律、単位律）の機械的検証
\item 型安全性による定義の明確化
\item モノトニシティなどの重要な性質の明示的証明
\end{itemize}

% ============================================================
\section{基本概念：ランク付き対象}
% ============================================================

\subsection{ランク付き対象の定義}

順序付き集合によるランク付けの概念を導入する。

\begin{definition}[前順序]
集合 $\alpha$ 上の二項関係 $\le$ が\emph{前順序}（preorder）であるとは、
以下を満たすことである：
\begin{itemize}
\item 反射律：任意の $a \in \alpha$ に対し $a \le a$、
\item 推移律：$a \le b$ かつ $b \le c$ ならば $a \le c$。
\end{itemize}
\end{definition}

\begin{definition}[ランク付き対象]\label{def:ranked}
前順序 $(\alpha, \le)$ に関する\emph{ランク付き対象}とは、対 $(X, r)$ である。
ここで：
\begin{itemize}
\item $X$ は型（集合）、
\item $r : X \to \alpha$ はランク関数。
\end{itemize}
Lean 4では、これを構造体 $\Ranked\ \alpha\ X$ として実装し、
フィールド $\mathtt{rank} : X \to \alpha$ を持つ。
\end{definition}

\begin{example}
$\alpha = \mathbb{N}$ を自然数の通常の順序とする。
$X = \{$有限グラフ$\}$ として、各グラフ $G$ のランクをその頂点数 $r(G) = |V(G)|$ と定義すれば、
$(X, r)$ はランク付き対象である。
\end{example}

\subsection{層の概念}

ランク関数から自然に導かれる重要な概念が層（layer）である。

\begin{definition}[層]\label{def:layer}
ランク付き対象 $R = (X, r)$ と $n \in \alpha$ に対し、
$n$-\emph{層} $\layer_R(n)$ を以下で定義する：
\[
  \layer_R(n) := \{\, x \in X \mid r(x) \le n \,\}.
\]
Lean記法では $\mathtt{R.layer}\ n$ と書く。
\end{definition}

\begin{remark}
層はランクの不等式を集合の帰属性に変換する。
これにより、存在層モルフィズム（後述の定義~\ref{def:rankhomd}）で
集合論的な「写像」の言葉で条件を述べることができる。
\end{remark}

\begin{proposition}[層のモノトニシティ]
$m \le n$ ならば $\layer_R(m) \subseteq \layer_R(n)$。
\end{proposition}

\begin{proof}
$x \in \layer_R(m)$ とすると $r(x) \le m$ である。
$m \le n$ より推移律から $r(x) \le n$ となり、$x \in \layer_R(n)$。
\end{proof}

% ============================================================
\section{データ層モルフィズム：\texorpdfstring{$\RankHomLe$}{RankHomLe}}
% ============================================================

\subsection{定義と基本性質}

3つの層の中で、$\RankHomLe$ を「真理の源泉」(source of truth) として採用する。

\begin{definition}[$\RankHomLe$ モルフィズム]\label{def:rankhomle}
前順序 $\alpha, \beta$ に関するランク付き対象 $R = (X, r_X)$ および $S = (Y, r_Y)$ の間の
\emph{データ層モルフィズム} $f : R \to S$ は、以下のデータからなる：
\begin{itemize}
\item 台写像 $f.\mathtt{map} : X \to Y$、
\item 添字写像 $f.\mathtt{indexMap} : \alpha \to \beta$、
\item モノトニシティ：$f.\mathtt{indexMap}$ は単調、
\item ランク不等式（緩可換性）：
\[
  r_Y(f.\mathtt{map}(x)) \le f.\mathtt{indexMap}(r_X(x))
  \qquad (\forall x \in X).
\]
\end{itemize}
\end{definition}

\begin{remark}
ランク不等式は、図式
\[
\begin{tikzcd}
X \arrow[r, "f.\mathtt{map}"] \arrow[d, "r_X"'] & Y \arrow[d, "r_Y"] \\
\alpha \arrow[r, "f.\mathtt{indexMap}"'] & \beta
\end{tikzcd}
\]
が「$\le$ の意味で可換」であることを意味する。
\end{remark}

\subsection{恒等射と合成}

圏の公理を満たすことを示すために、恒等射と合成を定義する。

\begin{definition}[恒等射]
ランク付き対象 $R = (X, r)$ に対し、恒等射 $\mathrm{id}_R : R \to R$ を以下で定義する：
\begin{itemize}
\item $(\mathrm{id}_R).\mathtt{map} = \mathrm{id}_X$、
\item $(\mathrm{id}_R).\mathtt{indexMap} = \mathrm{id}_\alpha$、
\item モノトニシティは自明（$a \le b \Rightarrow a \le b$）、
\item ランク不等式：$r(\mathrm{id}_X(x)) = r(x) \le r(x)$（反射律）。
\end{itemize}
\end{definition}

\begin{definition}[合成]\label{def:composition}
$f : R \to S$、$g : S \to T$ を $\RankHomLe$ モルフィズムとする。
合成 $g \circ f : R \to T$ を以下で定義する：
\begin{itemize}
\item $(g \circ f).\mathtt{map} = g.\mathtt{map} \circ f.\mathtt{map}$、
\item $(g \circ f).\mathtt{indexMap} = g.\mathtt{indexMap} \circ f.\mathtt{indexMap}$、
\item モノトニシティ：単調写像の合成は単調、
\item ランク不等式：任意の $x \in X$ に対し
\begin{align*}
r_T(g.\mathtt{map}(f.\mathtt{map}(x)))
&\le g.\mathtt{indexMap}(r_S(f.\mathtt{map}(x)))
  \quad (\text{$g$ のランク不等式})\\
&\le g.\mathtt{indexMap}(f.\mathtt{indexMap}(r_X(x)))
  \quad (\text{$f$ のランク不等式と $g.\mathtt{indexMap}$ のモノトニシティ}).
\end{align*}
\end{itemize}
\end{definition}

\begin{theorem}[合成の結合律]
$f : R \to S$、$g : S \to T$、$h : T \to U$ に対し、
$(h \circ g) \circ f = h \circ (g \circ f)$。
\end{theorem}

\begin{proof}
台写像と添字写像の等式は関数合成の結合律から従う。
モノトニシティとランク不等式は命題的内容（Prop）なので、証明無関係性により自動的に一致する。
Leanでは、構造体を分解（$\mathtt{cases}$）すれば定義的に等しい。
\end{proof}

\subsection{なぜ $\RankHomLe$ を選ぶのか}

$\RankHomLe$ を標準的な実装層として採用する理由は二つある：

\begin{enumerate}
\item \textbf{数学的安定性}：
合成の閉包性は、$\le$ の推移律と $\mathtt{indexMap}$ のモノトニシティのみで示される。
測度論的・解析的な追加構造は不要である。

\item \textbf{工学的安定性}：
合成の証明の核心は、定義~\ref{def:composition}で見たように
\[
  r_T(g(f(x))) \le g_*(r_S(f(x))) \le g_*(f_*(r_X(x)))
\]
という「一行の不等式連鎖」である。
これは外部補題にほとんど依存せず、ライブラリ更新に対して頑健である。
\end{enumerate}

% ============================================================
\section{圏 \texorpdfstring{$\RankCatLe$}{RankCatLe} の構成}
% ============================================================

\subsection{対象と射の定義}

前順序 $\alpha$ を固定する。

\begin{definition}[$\RankCatLe(\alpha)$ の対象]
圏 $\RankCatLe(\alpha)$ の対象は、依存対
\[
  \Obj(\RankCatLe(\alpha)) := \sum_{X : \Type} \Ranked(\alpha, X)
\]
である。すなわち、対象は台型 $X$ とランク関数 $r : X \to \alpha$ の対である。
\end{definition}

\begin{definition}[$\RankCatLe(\alpha)$ の射]
対象 $T = (X, r_X)$、$U = (Y, r_Y)$ に対し、
射の集合を
\[
  \Hom(T, U) := \RankHomLe(r_X, r_Y)
\]
と定義する。
\end{definition}

\subsection{圏構造の検証}

\begin{theorem}[$\RankCatLe$ は圏である]\label{thm:rankcatle-category}
$\RankCatLe(\alpha)$ は以下の構造を持つ圏である：
\begin{itemize}
\item 恒等射：各対象 $T$ に対し $\mathrm{id}_T \in \Hom(T, T)$、
\item 合成：$f \in \Hom(T, U)$、$g \in \Hom(U, V)$ に対し $g \circ f \in \Hom(T, V)$、
\item 単位律：$\mathrm{id}_U \circ f = f$ および $f \circ \mathrm{id}_T = f$、
\item 結合律：$(h \circ g) \circ f = h \circ (g \circ f)$。
\end{itemize}
\end{theorem}

\begin{proof}
前節で恒等射と合成を定義した。
単位律と結合律は、モルフィズムの構造体を分解すれば、
台写像と添字写像の等式は関数の単位律・結合律から従い、
残りのフィールド（$\mathtt{mono}$、$\mathtt{rank\_le}$）は命題なので自動的に一致する。
Lean 4では、$\mathtt{cases}$ タクティクと $\mathtt{rfl}$（反射律）で完了する。
\end{proof}

\begin{remark}
圏の公理の証明は、ランクの不等式の\emph{内部}を検査する必要がない。
これは、$\RankHomLe$ の設計が「正しい抽象レベル」にあることを示唆している。
\end{remark}

\subsection{忘却関手}

標準的な忘却関手を定義する。

\begin{definition}[台忘却関手]
関手 $U : \RankCatLe(\alpha) \to \Type$ を以下で定義する：
\begin{itemize}
\item 対象：$U(X, r) = X$、
\item 射：$U(f) = f.\mathtt{map}$。
\end{itemize}
\end{definition}

\begin{proposition}
$U$ は関手である。
\end{proposition}

\begin{proof}
$U(\mathrm{id}_T) = (\mathrm{id}_T).\mathtt{map} = \mathrm{id}_X = \mathrm{id}_{U(T)}$、
および $U(g \circ f) = (g \circ f).\mathtt{map} = g.\mathtt{map} \circ f.\mathtt{map} = U(g) \circ U(f)$。
\end{proof}

\begin{remark}
$U$ は一般には忠実（faithful）ではない。
なぜなら、異なる添字写像 $\mathtt{indexMap}$ を持つが同じ台写像を持つ射が存在しうるからである。
忠実な忘却関手を得るには、データフィールド $(\mathtt{map}, \mathtt{indexMap})$ の対を保持し、
証明フィールドのみを忘却すればよい。
\end{remark}

% ============================================================
\section{3つのモルフィズム層}
% ============================================================

異なる強度の条件を持つ3種類のモルフィズムを定義する。

\subsection{等式層：\texorpdfstring{$\RankHomEq$}{RankHomEq}}

最も強い条件を持つ層。

\begin{definition}[$\RankHomEq$ モルフィズム]\label{def:rankHomeq}
$\RankHomEq$ モルフィズム $f : R \to S$ は以下のデータからなる：
\begin{itemize}
\item 台写像 $f.\mathtt{map} : X \to Y$、
\item 添字写像 $f.\mathtt{indexMap} : \alpha \to \beta$、
\item モノトニシティ：$f.\mathtt{indexMap}$ は単調、
\item ランク等式（厳密可換性）：
\[
  r_Y(f.\mathtt{map}(x)) = f.\mathtt{indexMap}(r_X(x))
  \qquad (\forall x \in X).
\]
\end{itemize}
\end{definition}

\begin{remark}
$\RankHomEq$ では、ランクが添字写像によって\emph{等式}で保存される。
これは図式が厳密に可換であることを意味する。
\end{remark}

\subsection{不等式層：\texorpdfstring{$\RankHomLe$}{RankHomLe}}

中間の強度を持つ層（既に詳述）。

定義~\ref{def:rankhomle}で定義した $\RankHomLe$ は、
ランク不等式 $r_Y(f.\mathtt{map}(x)) \le f.\mathtt{indexMap}(r_X(x))$ を要求する。

\subsection{存在層：\texorpdfstring{$\RankHomD$}{RankHomD}}

最も弱い条件を持つ層。

\begin{definition}[$\RankHomD$ モルフィズム]\label{def:rankhomd}
$\RankHomD$ モルフィズム $f : R \to S$ は以下のデータからなる：
\begin{itemize}
\item 台写像 $f.\mathtt{map} : X \to Y$、
\item 層の存在条件：任意の $n \in \alpha$ に対し、
ある $m \in \beta$ が存在して、
\[
  \forall x \in X,\ r_X(x) \le n \Rightarrow r_Y(f.\mathtt{map}(x)) \le m.
\]
\end{itemize}
\end{definition}

\begin{remark}
$\RankHomD$ は添字写像を\emph{持たない}。
各ソースのランクレベル $n$ に対し、ターゲットのレベル $m$ が存在的に保証されるのみである。
\end{remark}

\begin{proposition}[層による特徴付け]
$\RankHomD$ の条件は、以下と同値である：
\[
  \forall n \in \alpha,\ \exists m \in \beta,\
  f.\mathtt{map}(\layer_R(n)) \subseteq \layer_S(m).
\]
すなわち、各層が何らかの層に写される。
\end{proposition}

\begin{proof}
定義~\ref{def:layer}より、
$x \in \layer_R(n)$ ⇔ $r_X(x) \le n$、
$y \in \layer_S(m)$ ⇔ $r_Y(y) \le m$。
よって条件は同値である。
\end{proof}

% ============================================================
\section{層の間のブリッジ}
% ============================================================

3つの層の間の自然な埋め込みを構成する。

\subsection{全体像}

\begin{figure}[h]
\centering
\begin{tikzcd}[column sep=3cm, row sep=1.5cm]
\RankHomEq \arrow[r, "\mathtt{toLe}"]
\arrow[rr, bend left=15, "\mathtt{toD}"above]
& \RankHomLe \arrow[r, "\mathtt{toD}"]
& \RankHomD
\end{tikzcd}
\caption{3つのモルフィズム層とブリッジ}
\label{fig:bridges}
\end{figure}

図~\ref{fig:bridges}に示すように、
等式層から不等式層へ、不等式層から存在層へと、
自然な「忘却」方向の写像が存在する。

\subsection{ブリッジ1：Eq → Le}

\begin{definition}[Eq → Le ブリッジ]\label{def:eq-to-le}
$f : R \to S$ を $\RankHomEq$ モルフィズムとする。
$f.\mathtt{toLe} : R \to S$ を以下の $\RankHomLe$ モルフィズムとして定義する：
\begin{itemize}
\item $(f.\mathtt{toLe}).\mathtt{map} = f.\mathtt{map}$、
\item $(f.\mathtt{toLe}).\mathtt{indexMap} = f.\mathtt{indexMap}$、
\item モノトニシティ：$f$ から継承、
\item ランク不等式：任意の $x$ に対し、
\[
  r_Y(f.\mathtt{map}(x))
  = f.\mathtt{indexMap}(r_X(x))
  \qquad (\text{$f$ のランク等式})
\]
より、反射律から
\[
  r_Y(f.\mathtt{map}(x)) \le f.\mathtt{indexMap}(r_X(x)).
\]
\end{itemize}
\end{definition}

\begin{remark}
等式は不等式を含意するため、$\mathtt{toLe}$ は自明なブリッジである。
Leanでは $\mathtt{le\_of\_eq}$ 補題を使用する。
\end{remark}

\subsection{ブリッジ2：Le → D}

不等式層から存在層への変換は、添字写像を証人（witness）として選ぶことで実現される。

\begin{definition}[Le → D ブリッジ]\label{def:le-to-d}
$f : R \to S$ を $\RankHomLe$ モルフィズムとする。
$f.\mathtt{toD} : R \to S$ を以下の $\RankHomD$ モルフィズムとして定義する：
\begin{itemize}
\item $(f.\mathtt{toD}).\mathtt{map} = f.\mathtt{map}$、
\item 層の存在条件：任意の $n \in \alpha$ に対し、
$m := f.\mathtt{indexMap}(n)$ を証人として選ぶ。
このとき、任意の $x \in X$ で $r_X(x) \le n$ ならば、
\begin{align*}
r_Y(f.\mathtt{map}(x))
&\le f.\mathtt{indexMap}(r_X(x))
  \qquad (\text{$f$ のランク不等式}) \\
&\le f.\mathtt{indexMap}(n)
  \qquad (\text{$f.\mathtt{indexMap}$ のモノトニシティと $r_X(x) \le n$}) \\
&= m.
\end{align*}
\end{itemize}
\end{definition}

\begin{lemma}[選ばれた証人]
$f : R \to S$ を $\RankHomLe$ モルフィズム、$n \in \alpha$ とする。
このとき、
\[
  \forall x \in X,\ r_X(x) \le n
  \Rightarrow r_Y(f.\mathtt{map}(x)) \le f.\mathtt{indexMap}(n).
\]
\end{lemma}

\begin{proof}
定義~\ref{def:le-to-d}の計算による。
\end{proof}

\subsection{合成ブリッジ：Eq → D}

便宜上、$\RankHomEq$ から $\RankHomD$ への直接的なブリッジも定義できる。

\begin{definition}[Eq → D ブリッジ]
$f : R \to S$ を $\RankHomEq$ モルフィズムとする。
$f.\mathtt{toD} := (f.\mathtt{toLe}).\mathtt{toD}$
と定義する。
\end{definition}

\begin{proposition}
図~\ref{fig:bridges}において、
$\RankHomEq$ から $\RankHomD$ への2つの経路は一致する：
\[
  f.\mathtt{toD} = (f.\mathtt{toLe}).\mathtt{toD}.
\]
\end{proposition}

\begin{proof}
定義より明らか。台写像はすべて $f.\mathtt{map}$ であり、
層の存在条件の証人の選び方も同一である。
\end{proof}

% ============================================================
\section{層の比較と設計原則}
% ============================================================

\subsection{強度の階層}

3つの層の間には、以下の含意関係がある：

\begin{theorem}[層の強度比較]
\[
  \RankHomEq \Rightarrow \RankHomLe \Rightarrow \RankHomD.
\]
すなわち、等式条件は不等式条件を含意し、不等式条件は存在条件を含意する。
\end{theorem}

\begin{proof}
ブリッジ $\mathtt{toLe}$ と $\mathtt{toD}$ の存在（定義~\ref{def:eq-to-le}、\ref{def:le-to-d}）による。
\end{proof}

\begin{remark}
逆方向の含意は一般には成り立たない。
例えば、$\RankHomD$ モルフィズムは添字写像を持たないため、
一般に $\RankHomLe$ に「持ち上げる」ことはできない。
\end{remark}

\subsection{各層の役割}

\begin{itemize}
\item \textbf{$\RankHomEq$（等式層）}：
ランクが厳密に保存される最も強い条件。
同型射や埋め込みの特徴付けに有用。

\item \textbf{$\RankHomLe$（不等式層）}：
「真理の源泉」として採用される中核的な層。
合成の閉包性が簡潔に示され、圏構造が自然に定まる。

\item \textbf{$\RankHomD$（存在層）}：
最も弱い層で、層を層に写すという集合論的条件のみを要求。
添字写像が自然に定義できない状況でも使用可能。
\end{itemize}

\subsection{T3時点での制限事項}

本研究（T3スナップショット）では、以下の内容を\emph{意図的に含めていない}：

\begin{itemize}
\item $\RankCatLe$ と既存の圏族（例：$\mathtt{Cat\_le}$、$\mathtt{Cat\_D}$、$\mathtt{Cat\_WithMin}$）との関係、
\item $\RankHomEq$、$\RankHomD$ に対する独立した圏構造の構築、
\item 各層の間の関手としての特徴付け。
\end{itemize}

これらは将来の作業項目（WBS: Work Breakdown Structure）として位置づけられている。
T3の目標は、クリーンなランク中心の圏とその内部的な層比較を確立することである。

% ============================================================
\section{形式化の詳細}
% ============================================================

\subsection{Lean 4における実装}

本稿の数学的内容は、すべてLean 4で形式化されている。
主なファイル構成は以下の通り：

\begin{description}
\item[\texttt{RankHomLe.lean}]
データ層モルフィズム $\RankHomLe$ の定義、恒等射、合成。

\item[\texttt{RankCatLe.lean}]
圏 $\RankCatLe$ の構築と圏公理の証明。

\item[\texttt{RankBridge\_EqLeD.lean}]
3つの層の定義とブリッジ関数 $\mathtt{toLe}$、$\mathtt{toD}$。

\item[\texttt{RankCat.lean}]
便利なラッパーとスケルトン（最小限の実装）。
\end{description}

\subsection{主要な証明のスケッチ}

\paragraph{合成の結合律（定理~\ref{thm:rankcatle-category}）}

Leanでの証明は以下のように進む：
\begin{verbatim}
assoc := by
  intro A B C D f g h
  cases f; cases g; cases h; rfl
\end{verbatim}

すなわち、モルフィズムを構造体として分解（$\mathtt{cases}$）し、
データフィールドの等式は関数合成の結合律から反射的に成り立ち（$\mathtt{rfl}$）、
証明フィールドは命題的に一意なので自動的に一致する。

\paragraph{ブリッジ Le → D の証明（定義~\ref{def:le-to-d}）}

Leanコードの核心部分：
\begin{verbatim}
def toD (f : RankHomLe R S) : RankHomD R S where
  map := f.map
  map_layer := by
    intro n
    refine ⟨f.indexMap n, ?_⟩
    intro x hx
    exact le_trans (f.rank_le x) (f.mono hx)
\end{verbatim}

証人として $f.\mathtt{indexMap}(n)$ を提示し、
不等式の連鎖 $r_Y(f.\mathtt{map}(x)) \le f.\mathtt{indexMap}(r_X(x)) \le f.\mathtt{indexMap}(n)$
を $\mathtt{le\_trans}$ で構成する。

\subsection{型安全性と依存型}

Lean 4の依存型システムにより、以下の安全性が保証される：

\begin{itemize}
\item ランク関数の値域の型（前順序 $\alpha$, $\beta$）が明示的にトラッキングされる。
\item モルフィズムの合成において、中間対象の型の整合性が型チェックで検証される。
\item 添字写像のモノトニシティが証明義務として明示的に要求される。
\end{itemize}

% ============================================================
\section{図式による可視化}
% ============================================================

\subsection{ランク可換図式}

$\RankHomLe$ モルフィズム $f : R \to S$ は、以下の図式の緩可換性を表す：

\begin{figure}[h]
\centering
\begin{tikzcd}[column sep=2.5cm, row sep=2cm]
X \arrow[r, "f.\mathtt{map}"] \arrow[d, "r_X"']
& Y \arrow[d, "r_Y"] \\
\alpha \arrow[r, "f.\mathtt{indexMap}"']
& \beta
\end{tikzcd}
\caption{$\RankHomLe$ の緩可換図式（$\le$ の意味で可換）}
\label{fig:rankhomle-diagram}
\end{figure}

ここで「$\le$ の意味で可換」とは、
$r_Y \circ f.\mathtt{map} \le f.\mathtt{indexMap} \circ r_X$
（点ごとの不等式）を意味する。

\subsection{合成の視覚化}

$f : R \to S$、$g : S \to T$ の合成を図示すると：

\begin{figure}[h]
\centering
\begin{tikzcd}[column sep=2cm, row sep=1.8cm]
X \arrow[r, "f.\mathtt{map}"] \arrow[d, "r_X"']
& Y \arrow[r, "g.\mathtt{map}"] \arrow[d, "r_Y"]
& Z \arrow[d, "r_Z"] \\
\alpha \arrow[r, "f.\mathtt{indexMap}"']
& \beta \arrow[r, "g.\mathtt{indexMap}"']
& \gamma
\end{tikzcd}
\caption{合成 $g \circ f$ の図式}
\label{fig:composition}
\end{figure}

外側の長方形も $\le$ の意味で可換となる。

\subsection{層の包含関係}

層のブリッジを図示すると：

\begin{figure}[h]
\centering
\begin{tikzpicture}[node distance=2.5cm and 3.5cm]
  \node (eq) [draw, rectangle, minimum width=3cm, minimum height=1.2cm, align=center]
    {\textbf{等式層} \\ $\RankHomEq$ \\ $r_Y(f(x)) = \varphi(r_X(x))$};
  \node (le) [draw, rectangle, minimum width=3cm, minimum height=1.2cm, align=center, below=of eq]
    {\textbf{不等式層} \\ $\RankHomLe$ \\ $r_Y(f(x)) \le \varphi(r_X(x))$};
  \node (d) [draw, rectangle, minimum width=3cm, minimum height=1.2cm, align=center, below=of le]
    {\textbf{存在層} \\ $\RankHomD$ \\ $\forall n\,\exists m\,\forall x\ldots$};

  \draw[->, thick, >=stealth] (eq) -- (le) node[midway, right] {$\mathtt{toLe}$};
  \draw[->, thick, >=stealth] (le) -- (d) node[midway, right] {$\mathtt{toD}$};
  \draw[->, thick, >=stealth, dashed] (eq) to[bend right=30] (d) node[midway, left] {$\mathtt{toD}$};

  \node[right=0.5cm of eq] {\footnotesize 最強};
  \node[right=0.5cm of d] {\footnotesize 最弱};
\end{tikzpicture}
\caption{3つの層の強度階層}
\label{fig:layer-hierarchy}
\end{figure}

% ============================================================
\section{応用と今後の展望}
% ============================================================

\subsection{計算複雑性理論への応用}

ランク付き構造は、計算複雑性のクラス間の還元を記述するのに適している。
例えば、問題 $A$ から問題 $B$ への多項式時間還元を、
ランク $=$ 計算時間として $\RankHomLe$ モルフィズムでモデル化できる可能性がある。

\subsection{順序構造の圏論}

$\RankCatLe$ は、順序集合の圏（posets）と通常の圏論の中間に位置する構造を提供する。
今後の研究では、以下の方向が考えられる：

\begin{itemize}
\item 極限・余極限の存在性、
\item モナド構造の導入、
\item エンリッチ圏としての特徴付け。
\end{itemize}

\subsection{他の圏族との統合}

既存の圏族（$\mathtt{Cat\_le}$、$\mathtt{Cat\_WithMin}$ など）と $\RankCatLe$ の関係を明確化し、
忘却関手や比較関手を構成することが次のステップである。

\subsection{形式化の拡張}

Lean 4での形式化を拡張し、以下を目指す：

\begin{itemize}
\item $\RankHomEq$、$\RankHomD$ の圏構造、
\item ブリッジを関手として定式化、
\item 随伴関手の構成（もし存在すれば）。
\end{itemize}

% ============================================================
\section{結論}
% ============================================================

本稿では、ランク付き構造の圏論的枠組みを構築した。
主な成果は以下の通りである：

\begin{enumerate}
\item \textbf{データ層モルフィズム $\RankHomLe$ の定義}：
ランク不等式による緩可換性を要求し、簡潔な合成証明を実現。

\item \textbf{圏 $\RankCatLe$ の構成}：
圏の公理をLean 4で機械的に検証し、構造の健全性を確立。

\item \textbf{3つの層の定義とブリッジ}：
等式層 $\RankHomEq$、不等式層 $\RankHomLe$、存在層 $\RankHomD$ の間の
自然な埋め込みを構成し、強度の階層を明確化。
\end{enumerate}

本研究は、形式化数学の利点を活用し、
抽象的な圏論的構造を具体的な計算オブジェクト（Leanコード）として実現した。
今後の展開により、計算複雑性理論や順序構造の研究への応用が期待される。

% ============================================================
% 参考文献（簡略版）
% ============================================================

\begin{thebibliography}{9}

\bibitem{lean4}
Leonardo de Moura, Sebastian Ullrich,
\textit{The Lean 4 Theorem Prover and Programming Language},
\url{https://leanprover.github.io/}, 2021.

\bibitem{mathlib}
The mathlib Community,
\textit{The Lean Mathematical Library},
\url{https://github.com/leanprover-community/mathlib4}, 2024.

\bibitem{maclane}
Saunders Mac Lane,
\textit{Categories for the Working Mathematician},
Springer, 2nd edition, 1998.

\bibitem{categorytheory-lean}
Scott Morrison, et al.,
\textit{Category Theory in Lean},
Mathlib documentation, 2024.

\end{thebibliography}

% ============================================================
% 付録
% ============================================================

\appendix

\section{Leanコードの抜粋}

以下に、主要な定義のLeanコードを示す。

\subsection{\texttt{RankHomLe}の定義}

\begin{verbatim}
structure RankHomLe {α : Type v} {β : Type w}
    {X : Type u} {Y : Type u}
    [Preorder α] [Preorder β]
    (R : Ranked α X) (S : Ranked β Y) where
  map      : X → Y
  indexMap : α → β
  mono     : Monotone indexMap
  rank_le  : ∀ x : X, S.rank (map x) ≤ indexMap (R.rank x)
\end{verbatim}

\subsection{合成の定義}

\begin{verbatim}
def comp {R : Ranked α X} {S : Ranked β Y} {T : Ranked γ Z}
    (f : RankHomLe R S) (g : RankHomLe S T) :
    RankHomLe R T where
  map := g.map ∘ f.map
  indexMap := g.indexMap ∘ f.indexMap
  mono := g.mono.comp f.mono
  rank_le := by
    intro x
    have hf : S.rank (f.map x) ≤ f.indexMap (R.rank x) :=
      f.rank_le x
    have hg : T.rank (g.map (f.map x)) ≤
      g.indexMap (S.rank (f.map x)) := g.rank_le (f.map x)
    exact le_trans hg (g.mono hf)
\end{verbatim}

\subsection{ブリッジ Le → D}

\begin{verbatim}
def toD (f : RankHomLe R S) : RankHomD R S where
  map := f.map
  map_layer := by
    intro n
    refine ⟨f.indexMap n, ?_⟩
    intro x hx
    exact le_trans (f.rank_le x) (f.mono hx)
\end{verbatim}

\section{記号表}

\begin{description}
\item[$\alpha, \beta, \gamma$] 前順序型（ランクの値域）
\item[$X, Y, Z$] 台型（carrier type）
\item[$r, r_X, r_Y$] ランク関数
\item[$\Ranked(\alpha, X)$] 型 $X$ 上の $\alpha$-ランク付き構造
\item[$\layer_R(n)$] ランク付き対象 $R$ の $n$-層
\item[$\RankHomLe$] データ層モルフィズム（不等式による）
\item[$\RankHomEq$] 等式層モルフィズム
\item[$\RankHomD$] 存在層モルフィズム
\item[$\RankCatLe(\alpha)$] $\alpha$-ランク付き対象の圏
\item[$\mathtt{map}$] 台写像
\item[$\mathtt{indexMap}$] 添字写像
\item[$\mathtt{toLe}$] Eq → Le ブリッジ
\item[$\mathtt{toD}$] Le → D または Eq → D ブリッジ
\end{description}

\end{document}
